% Source: https://artofproblemsolving.com/community/c3h2285159_equilateral_triangle_problem
% Stable identity: aops-topic-2285159-q01
Let the horizontal side be $PR$ with $P=(p,q)$ and $R=(p+2a,q)$, $a>0$, and let $T$ be the third vertex. Since $TP=TR$, the vertex $T$ lies on the perpendicular bisector $x=p+a$ of $PR$. The height of an equilateral triangle of side $2a$ is $\sqrt{(2a)^2-a^2}=a\sqrt3$, so $T=(p+a,\ q\pm a\sqrt3)$.

\begin{asy}
usepackage("amsmath");
size(6cm);
defaultpen(fontsize(10pt));
pair P=(0.6,0.5), R=(2.6,0.5), T=(1.6,0.5+sqrt(3));
draw((-0.3,0)--(3.4,0),Arrow(6)); draw((0,-0.3)--(0,2.7),Arrow(6));
label("$x$",(3.4,0),E); label("$y$",(0,2.7),N);
draw(P--R--T--cycle,linewidth(0.9));
draw((1.6,0.5)--T,dashed);
label("$P(p,q)$",P,SW); label("$R(p+2a,q)$",R,SE); label("$T(p+a,q+a\sqrt3)$",T,N);
draw(arc(P,0.3,0,60)); label("$60^\circ$",P+0.5*dir(30),fontsize(9pt));
label("$a$",(1.1,0.5),S); label("$a\sqrt3$",(1.6,1.3),E,fontsize(9pt));
\end{asy}

\textbf{(b)} If $T$ is above $PR$, the slopes are
\[\text{slope }PT=\dfrac{a\sqrt3}{a}=\sqrt3,\qquad \text{slope }RT=\dfrac{a\sqrt3}{-a}=-\sqrt3;\]
if $T$ is below, they are $-\sqrt3$ and $\sqrt3$. In every case the other two sides have slopes $\boxed{\sqrt3\text{ and }-\sqrt3}$. (This also follows from the $60^\circ$ angles at $P$ and $R$: the slopes are $\pm\tan60^\circ$.)

\textbf{(a)} Since $\sqrt3\ne1$, the claim is false; for example, the triangle with vertices $(0,0)$, $(2,0)$, $(1,\sqrt3)$ is equilateral with a horizontal side, and its other sides have slopes $\sqrt3$ and $-\sqrt3$. Slopes $1$ and $-1$ would make $45^\circ$ angles with the horizontal side, and the triangle would be the right isosceles triangle with angles $45^\circ,45^\circ,90^\circ$, not an equilateral one.
